\section*{KI-Verzeichnis}

Für die Arbeit „Walking breaks \& 100\% focus\_time“ wurde das generative KI-Modell Claude (Anthropic) über Claude Code mit dem Plugin CoWriting ausschließlich assistierend eingesetzt: für Hinweise, Rückfragen, Erklärungen und Literaturhilfe. Der Text des Manuskripts wurde von der Autorin bzw. dem Autor verfasst. Das Plugin weist Schreibzugriffe der KI auf Manuskriptdateien ab, protokolliert jede Änderung, die es während eines KI-Werkzeugaufrufs dennoch feststellt (siehe unten), und die Versionsgeschichte des Manuskripts (git) zeigt, wie der Text entstanden ist. Alle KI-Eingaben wurden automatisch protokolliert (ai-log/).

\begin{tabular}{p{.22\linewidth}p{.18\linewidth}p{.2\linewidth}p{.2\linewidth}p{.1\linewidth}}
\hline
\textbf{Zweck} & \textbf{Werkzeug / Modell} & \textbf{Umfang} & \textbf{Verwendung} & \textbf{Kapitel} \\
\hline
Prüfung der Fachterminologie & Claude Code (seed) & 1 Anfrage, 2026-11-01 bis 2026-11-01 & nicht annotiert: 1 & – \\
\hline
\end{tabular}

\noindent(1) Reine Sprachkorrekturen (Grammatik/Rechtschreibung) müssen laut DFG in Anträgen nicht offengelegt werden; manche Hochschulen verlangen es dennoch.

\subsection*{Beispielhafte Anfragen}
\begin{itemize}
  \item \textit{Prüfung der Fachterminologie:} For this journal I always write participants, never subjects. Please keep that in mind for this paper.
\end{itemize}
\subsection*{Von der KI formulierte Textstellen ([AI-TEXT])}
Keine.
\subsection*{Während KI-Werkzeugaufrufen festgestellte Änderungen am Manuskript}
Keine.
\subsection*{Nach einem Werkzeugaufruf rückgängig gemachte Änderungen am KI-Protokoll}
Keine.
\subsection*{Vom Schutz abgewiesene Schreibversuche der KI}
Keine.

\textit{Die Verantwortung für Inhalt, Quellen und Formulierungen liegt vollständig bei der Autorin bzw. dem Autor. Bitte an die Vorgaben der Zielinstitution anpassen.}
